\documentclass[11pt]{article}

\usepackage{hw_style}

% Set due date for this homework (updated to 2025)
\setduedate{11}{26}

% Setup homework number and header
\hwnumber{7}

\begin{document}

\begin{hwproblem}{1}
    计算下列表达式的拉普拉斯变换的像函数
    \begin{enumerate}
      \item $$
      \frac{1}{t}\left(e^{\beta t}-e^{\alpha t}\right) ;
      $$
      \item $$
      \int_0^t \frac{\sin u}{u} d u;
      $$
      \item 
      $$
      \frac{1}{\sqrt{\pi t}}.
      $$
    \end{enumerate}
\end{hwproblem}

\vspace{0.5cm}

\begin{hwproblem}{2}
    计算下列表达式的拉普拉斯变换的原函数
    \begin{enumerate}
      \item $$
      \bar{f}(p)=\frac{3 p}{p^2-1} ;
      $$
      \item $$
      \frac{4 p-1}{\left(p^2+p\right)\left(4 p^2-1\right)}.
      $$
    \end{enumerate}
\end{hwproblem}

\vspace{0.5cm}

\begin{hwproblem}{3}
    利用拉普拉斯变换计算:  $$\int_0^{\infty} \frac{\sin x t}{t} dt. $$
\end{hwproblem}

\vspace{0.5cm}
   
\begin{hwproblem}{4}
用拉普拉斯反变换中的Bromwich(又称黎曼-梅林)积分求
$$
\frac{p}{p^2 - \omega^2}
$$
的原函数.
\end{hwproblem}

\vspace{0.5cm}

\begin{hwproblem}{5}
    利用拉普拉斯变换求解微分方程
    $$
      \frac{d^3 y}{d t^3}+3 \frac{d^2 y}{d t^2}+3 \frac{d y}{d t}+y=6 e^{-t}, 
      \quad y(0)=\left.\frac{d y}{d t}\right|_{t=0}=\left.\frac{d^2 y}{d t^2}\right|_{t=0}=0 .
    $$
\end{hwproblem}

\vspace{0.5cm}
  

\begin{hwproblem}{6}
    将 $\delta(x)$ 展为实数形式的傅里叶积分.
\end{hwproblem}

\vspace{0.5cm}

\begin{hwproblem}{7}
    用$\Gamma$函数求积分
    $$
    \int_0^\infty x^{\alpha -1 } e^{-x \cos{\theta}} \cos\left( x \sin{\theta} \right) dx,
    $$
    其中$\alpha > 0, -\frac{\pi}{2} < \theta < \frac{\pi}{2}$.  
\end{hwproblem}

\vspace{0.5cm}

  
  
\begin{hwproblem}{8}
用$B$函数求积分
$$
\int_0^{\frac{\pi}{2}} \tan^\alpha{\theta} \,  d\theta, \quad  -1 < \alpha < 1.
$$
\end{hwproblem}

\vspace{0.5cm}
    

% \begin{hwproblem}{9}
% \end{hwproblem}

\vspace{0.5cm}

% \begin{hwproblem}{10}
% \end{hwproblem}

\vspace{0.5cm}

% \begin{hwproblem}{11}
% \end{hwproblem}

\vspace{0.5cm}

\hwrequirements

\end{document}